% META: {"id": "TSPE-GEOESPACE-ME-007", "chapitre": "TSPE-GEOMETRIE-ESPACE", "type_objet": "methode", "capacites_codes": ["C10"], "methodes": ["M7"], "sources_inspiration": [], "status": "generated"}
% BEGIN-VERIFY
% from sympy import Matrix
% i=Matrix([1,0,0]); j=Matrix([0,1,0]); v=Matrix([0,0,5]); u=Matrix([0,1,1])
% assert v.dot(i) == 0 and v.dot(j) == 0        # v orthogonal au plan z=0
% assert Matrix.hstack(v, Matrix([0,0,1])).rank() == 1
% assert u.dot(i) == 0 and u.dot(j) != 0        # le piege : une seule droite
% END-VERIFY
\begin{fichemethode}{M7}{Démontrer l'orthogonalité d'une droite et d'un plan}

\quandUtiliser{%
  Pour établir qu'une droite est orthogonale à un plan, ou pour exploiter cette
  orthogonalité dans un problème de configuration ou de lieu géométrique.
}

\pasApas{%
  \begin{enumerate}
    \item \textbf{Déterminer} un vecteur directeur $\vec{u}$ de la droite.
    \item \textbf{Choisir deux vecteurs non colinéaires} du plan, ou lire
          directement son vecteur normal si le plan est donné par une équation.
    \item \textbf{Calculer les deux produits scalaires} de $\vec{u}$ avec ces
          deux vecteurs.
    \item \textbf{Conclure} seulement si les \emph{deux} sont nuls~: la droite
          est alors orthogonale au plan, et $\vec{u}$ en est un vecteur normal.
    \item \textbf{Variante} si le plan est donné par $ax+by+cz+d=0$~: il suffit
          de vérifier que $\vec{u}$ est colinéaire à $\vec{n}(a;b;c)$.
  \end{enumerate}
}

\verifier{%
  Contrôler que les deux vecteurs du plan choisis sont bien non colinéaires~:
  deux vecteurs colinéaires ne donnent qu'une seule direction et la conclusion
  serait fausse. La variante par le vecteur normal fournit un second contrôle
  indépendant.
}

\pieges{%
  Être orthogonal à \emph{une} droite du plan ne suffit pas. Le plan possède
  une infinité de directions~: $\vec{u}(0;1;1)$ est orthogonal à l'axe des
  abscisses sans l'être au plan $z=0$, puisque $\vec{u}\cdot(0;1;0)=1$.
}

\exempleRedige{%
  Soit $\mathcal{P}$ le plan $z=0$, dirigé par $\vec{i}(1;0;0)$ et
  $\vec{j}(0;1;0)$, et $\vec{v}(0;0;5)$.
  \[
    \vec{v}\cdot\vec{i}=0,\qquad \vec{v}\cdot\vec{j}=0 .
  \]
  Les deux produits sont nuls et $\vec{i}$, $\vec{j}$ ne sont pas colinéaires~:
  la droite dirigée par $\vec{v}$ est orthogonale à $\mathcal{P}$. On le
  retrouve par la variante~: $\vec{v}=5\,\vec{n}$ avec $\vec{n}(0;0;1)$.
}

\end{fichemethode}
